\chapter{Discussion}
\label{chap:discussion}

This chapter interprets the results presented in Chapter~\ref{chap:results}
across the ten benchmark geometries and four experimental modes.
We follow the same domain order as the results chapter.
For each group we discuss what the numbers show, how the modes compare, and
what physical or architectural properties explain the findings.
The final section draws cross-domain conclusions.

% ─────────────────────────────────────────────────────────────────────────────
\section{2D NAS Architecture Transfer and Adaptive-$k$}
\label{sec:disc_2d_adaptive}

\subsection{What the results show}

We evaluated the adaptive-$k$ strategy on four 2D domains using four architecture
configurations: the NAS-PINN transfer architecture from Wang et
al.~\cite{wang2023} and our three NAS-found architectures.
All four achieved 83--85\% FEM saving with mean-window MAE below
$2.5\,^\circ\mathrm{C}$ across all domains.

We observed a geometry-dependent pattern.
On the Square and Circle domains, all four configurations produced nearly
identical results.
MAE differences were below $0.2\,^\circ\mathrm{C}$.
We found that the NAS transfer architecture generalised well to these smooth geometries.

On the L-Shape domain, Bayesian/TPE produced the lowest MAE across both
adaptive-$k$ ($1.36\,^\circ\mathrm{C}$, ten-seed mean) and single-seed runs.
The NSGA-II and NSGA-III architectures reached $2.34$--$2.44\,^\circ\mathrm{C}$
on the same domain.
We observed the same pattern on the Flower domain.

\subsection{Why adaptive-$k$ outperforms fixed-$k$ here}

We confirmed the advantage of adaptive-$k$ using the fixed-$k$ results
on the same domains.
At $k=1$, all architectures stay below $2.5\,^\circ\mathrm{C}$.
At $k=2$ and above, errors rise steeply.
At $k=5$, some configurations exceed $35\,^\circ\mathrm{C}$.
This is completely outside any useful range.

The physical reason is the temporal structure of the quench.
These surfaces are geometrically simple and cool rapidly.
The temperature gradient is steepest in the first few seconds.
A fixed $k=5$ schedule forces the PINN to predict over a 7.5-second window
even during this steep early phase.
The PINN drifts from the FEM reference over such a long unsupported interval.
The adaptive controller avoids this.
It holds $k=1$ during the steep gradient and promotes to $k=4$--$5$ only once
the field relaxes.
All four architectures converged to the same schedule on every domain.
We concluded that the controller responds to the physics, not to architecture-specific artefacts.

\subsection{Why Bayesian/TPE leads on the L-Shape}

The L-Shape contains a 270° reentrant corner.
This corner creates a concentrated, localised temperature gradient.
The large ReLU network of Bayesian/TPE learns this concentrated structure
effectively.
The compact tanh networks of NSGA-II and NSGA-III lack the capacity to
represent it with the same precision.
The NAS validation domain is a smooth rectangle.
It does not share this feature.
This explains why the transfer architecture, also selected on the rectangle,
underperforms Bayesian/TPE on the L-Shape.
We found that geometry-specific NAS adds value when the target geometry
has features absent from the validation domain.

% ─────────────────────────────────────────────────────────────────────────────
\section{Canonical 2D Benchmark Results}
\label{sec:disc_2d}

\subsection{What the results show}

The systematic fixed-$k$ sweep on the Rectangle, Circle, and L-Shape domains
confirmed the pattern identified above.
At $k=1$, all architectures achieve low errors, and Bayesian/TPE gives the
lowest mean MAE on all three canonical 2D domains ($2.12$--$2.75\,^\circ\mathrm{C}$).
Through $k=3$, Bayesian/TPE remains below $5\,^\circ\mathrm{C}$ on Rectangle,
Circle, and L-Shape, corresponding to a robust 65\% FEM saving.

As $k$ increased, we observed a clear divergence between architectures.
Bayesian/TPE degrades slowly on Rectangle and L-Shape, reaching
$4.03\,^\circ\mathrm{C}$ and $5.41\,^\circ\mathrm{C}$ at $k=5$.
The Circle is more sensitive to long unsupported windows: Bayesian/TPE rises
to $10.01\,^\circ\mathrm{C}$ at $k=5$.
The compact NSGA-II and NSGA-III architectures show larger degradation on the
harder cases, especially L-Shape, where NSGA-II reaches
$14.37\,^\circ\mathrm{C}$ at $k=5$.

\subsection{Why compact architectures degrade faster}

The canonical 2D sweep shows that Bayesian/TPE is the most robust architecture
across increasing skip factors.
Its larger ReLU network can keep the error controlled over longer windows,
whereas the compact tanh networks lose accuracy more quickly on geometries
with concentrated gradients.
At longer windows, the temperature field evolves in ways that require broader
spatial representation.
The smaller networks cannot track this evolution with the same margin.
The tanh activation can also saturate at large inputs, limiting each neuron's
effective range.

The L-Shape confirmed this interpretation.
The concentrated corner gradient is a clear and learnable target for a
high-capacity network with FEM anchoring.
NSGA-II, however, reaches $14.37\,^\circ\mathrm{C}$ at $k=5$.
This is the worst mean among the canonical 2D configurations.
The compact architecture cannot maintain accuracy when the corner gradient
must be predicted over a long unsupported window.

% ─────────────────────────────────────────────────────────────────────────────
\section{Canonical 3D Benchmark Results}
\label{sec:disc_3d}

\subsection{What the results show}

The 3D canonical benchmarks scale the problem to three spatial dimensions.
More collocation points, more complex boundary geometries, and a harder
training problem emerge within the same epoch budget.

At $k=1$, we measured 3D errors ranging from $4.41$ to $10.51\,^\circ\mathrm{C}$
(ten-seed means), compared to $1.64$--$3.52\,^\circ\mathrm{C}$ in 2D.
We found that Bayesian/TPE reached $k=5$ (80\% FEM saving) on all four 3D domains.
NSGA-II and NSGA-III were reliable through $k=3$--$4$ on most domains.

Two domains produced particularly notable results.
On the 3D L-Shape Prism, NSGA-III achieved the lowest $k=1$ error
($4.41\,^\circ\mathrm{C}$), outperforming Bayesian/TPE ($10.20\,^\circ\mathrm{C}$).
This is the only canonical domain where Bayesian/TPE is not the best at $k=1$.
On Cylinder and Stacked Cubes, NSGA-II showed high inter-seed variance and
above-threshold means at $k\geq4$.

The adaptive-$k$ controller applied to all four 3D domains with Bayesian/TPE
produced physically consistent schedules.
It held $k=1$ during the steep early transient and promoted to longer windows
as the field relaxed.

\subsection{Why the 3D L-Shape is different}

We observed that the 3D L-Shape prism concentrates its dominant gradient along a single
reentrant edge.
This is a spatially localised feature.
The smooth tanh basis of NSGA-III represents localised features compactly.
The 2D rectangle is the NAS validation domain.
It also has its dominant gradient at the boundary with a near-uniform interior.
NSGA-III's inductive bias therefore matches both the validation domain and the
3D L-Shape.

At $k=5$, however, NSGA-II's ten-seed mean reached $15.65\,^\circ\mathrm{C}$
($\sigma=8.85\,^\circ\mathrm{C}$) compared to Bayesian/TPE's
$13.96\,^\circ\mathrm{C}$.
The advantage of NSGA-III at $k=1$ does not transfer to higher skip factors.
Architecture selection must account for the target skip factor, not only the
$k=1$ baseline.

\subsection{NSGA-II failure on Cylinder and Stacked Cubes}

NSGA-II showed high inter-seed variance on the Cylinder ($\sigma$ rising to
$12.28\,^\circ\mathrm{C}$ at $k=5$) and above-threshold means on Stacked Cubes.
Two structural factors interact.

First, these geometries produce gradients that vary simultaneously in three
spatial directions: curved cylindrical surfaces and internal material interfaces.
The tanh activation saturates at large inputs.
This limits each neuron's contribution to steep, spatially concentrated gradients.
Few layers are insufficient to represent such structure.

Second, the compact architecture (${\approx}48$K parameters) produces a loss
landscape with many local minima of similar depth.
Different random initialisations converge to qualitatively different solutions.
The high inter-seed variance is the clearest evidence of this.

Bayesian/TPE avoids both problems.
ReLU does not saturate for positive inputs.
The larger parameter count (${\approx}93$K) gives the optimiser a smoother
trajectory.
The 3D L-Shape prism, where NSGA-II performs comparably to other architectures,
lacks the curved normals and material interfaces that trigger this failure.

\subsection{The degradation ratio across 3D domains}

We noted an important observation when comparing $k=1$ and $k=5$ errors in 3D.
On the Rectangular Prism, Bayesian/TPE's mean error moves from $8.42$ to
$9.88\,^\circ\mathrm{C}$, a ratio of only $1.2$.
On the Cylinder, the ratio is $1.8$.
These ratios are much smaller than the corresponding 2D values for compact
architectures ($4.8$--$7.6$).
Complex 3D geometries with rich spatial structure allow the PINN to represent
the field accurately even over longer windows, provided the architecture has
sufficient capacity.

% ─────────────────────────────────────────────────────────────────────────────
\section{Thermal Fin Benchmark}
\label{sec:disc_thermalfin}

\subsection{Fixed-$k$ results}

The key finding is that Bayesian/TPE with Fourier embedding is the only
architecture that stays within the $15\,^\circ\mathrm{C}$ threshold at all
five skip factors, while NSGA-II and NSGA-III top out at $k=4$
(Table~\ref{tab:tf_consolidated}).
What the summary table cannot show is the per-window error profile, which
explains why the Thermal Fin is harder than any canonical domain.

We observed that the first prediction window consistently produced the highest error across all
architectures: $17.7\,^\circ\mathrm{C}$ for Bayesian/TPE, $21.0\,^\circ\mathrm{C}$
for NSGA-II, and $27.1\,^\circ\mathrm{C}$ for NSGA-III.
The fin-tip gradient is steepest at quench onset, when the surface cools
instantly while the interior remains at $540\,^\circ\mathrm{C}$.
After the first few windows, all three architectures settled to steadily
decreasing errors.
By the final window they reached $6.35$, $4.50$, and $6.06\,^\circ\mathrm{C}$.
The mean-window MAE therefore averages a difficult early phase and an easier
late phase.

The Fourier embedding reduces Bayesian/TPE error by $1.8$--$4.4\,^\circ\mathrm{C}$
across all skip factors (Appendix~\ref{app:ablation}).
The same embedding is detrimental for compact tanh networks because the
128-dimensional sinusoidal projection overwhelms the three-layer capacity.

\subsection{Adaptive-$k$ results}

We achieved 45\% FEM saving with all three architectures with mean-window MAE within the
$15\,^\circ\mathrm{C}$ threshold.
The controller held $k=1$ through the first five or six windows and promoted
to $k=2$, $k=3$, then $k=4$ as the field relaxed.
NSGA-III needed one extra window at $k=1$ because its first-window error
($27.1\,^\circ\mathrm{C}$) was the largest and took longer to fall below the
promotion threshold.

We found that the adaptive schedule achieved 45\% FEM saving while the best
fixed-$k$ configuration achieves 80\%.
The root cause is a threshold--architecture mismatch.
A single promotion threshold was set for all architectures.
The three networks operate at different baseline error levels.
Bayesian/TPE's per-window error at $k=1$ sits near the threshold and promotes
early.
NSGA-III's error is consistently higher and triggers demotion where
Bayesian/TPE would promote.
An architecture-aware threshold that initialises from the known fixed-$k$ optimum
would decouple these effects.

\subsection{PINN-only results}

Removing FEM corrections exposed a critical limitation of the framework.
We measured the best mean-window result without anchoring as $59.20\,^\circ\mathrm{C}$
for NSGA-III at $k=5$, which is nearly four times the acceptance threshold.
Bayesian/TPE produced errors above $200\,^\circ\mathrm{C}$ at all skip factors.

The $5{\times}151$ ReLU network that performs best with FEM corrections
accumulated drift the fastest without them.
Each window's residual compounds into the next rather than being reset by the
FEM anchor.
The Thermal Fin PINN-only results confirm what the 2D preliminary run already
suggested (Section~\ref{subsec:2d_adap_pinn_only}): FEM anchoring is a
structural requirement of the pipeline, not an optional component.

% ─────────────────────────────────────────────────────────────────────────────
\section{Cross-Domain Analysis and MSWP Significance}
\label{sec:disc_nas_transfer}

\subsection{Architecture transfer across all domains}

We selected all three architectures on the 2D rectangle domain at $k=1$ and
deployed them across all eight benchmarks without re-running NAS.
We found that Bayesian/TPE achieved the lowest ten-seed mean MAE in 7 of 8 cases.
It was within $6\,^\circ\mathrm{C}$ of the best in the remaining one
(3D L-Shape at $k=1$, where NSGA-III led).
A high-capacity ReLU network selected on a simple geometry generalises well
to more complex ones because its expressivity is not the binding constraint on
any benchmark.

NSGA-II and NSGA-III transferred less reliably.
Both were selected by minimising parameter count as a secondary objective.
This penalises exactly the capacity needed to represent complex 3D gradient
fields.

\subsection{From simple to complex: the role of MSWP}

Across all ten benchmark geometries, we identified a consistent pattern as geometric complexity increases.

On simple 2D domains, we found that fixed-$k$ at high values fails dramatically.
The $k=5$-to-$k=1$ error ratio exceeds 4--7 for compact architectures.
Adaptive-$k$ resolves this by matching window length to the local gradient
magnitude.

On complex 3D domains and the Thermal Fin, the same ratio for Bayesian/TPE
falls to 1.2--2.5.
The richer spatial structure allows the PINN to predict accurately over longer
intervals.
At the same time, each FEM anchor call on an industrial mesh is far more
expensive.
The combination of manageable error growth and high per-call cost makes MSWP
most beneficial precisely where the geometry is complex.

\subsection{Practical guidance}

Based on the results, we draw three conclusions for practitioners.

First, Bayesian/TPE is the recommended architecture for all deployment scenarios.
Its accuracy--efficiency frontier dominates across the full range of geometries
and skip factors.

Second, adaptive-$k$ is the recommended scheduling strategy for simple and
intermediate geometries.
It combines the accuracy of $k=1$ during the steep early transient with the
FEM savings of higher $k$ during field relaxation.
For complex geometries, fixed $k=4$--$5$ with Bayesian/TPE is also viable and
achieves higher FEM savings.

Third, periodic FEM anchoring is essential.
The PINN-only results demonstrate that removing FEM corrections causes
catastrophic error accumulation on all tested geometries.
The framework is a FEM-reduction strategy, not a FEM-replacement strategy.
